\chapter{线性系统的根轨迹法}

% ===== 局部排版优化（仅影响本章，使版面更舒展）=====
\setlength{\parskip}{0.8ex plus 0.3ex minus 0.2ex}
\setlist{leftmargin=2.5em, itemsep=0.6em, parsep=0.3em, topsep=0.8em, partopsep=0.3em}
\renewcommand{\arraystretch}{1.5}
\setlength{\tabcolsep}{10pt}
\addtolength{\abovedisplayskip}{2pt}
\addtolength{\belowdisplayskip}{2pt}
% ========================================

\bigskip

\section{本章考核要点}
\begin{enumerate}
    \item 根轨迹的定义、根轨迹方程，以及相角条件与模值条件各自的地位；
    \item 闭环零、极点与开环零、极点之间的关系；
    \item 绘制 $180^\circ$ 常规根轨迹的八条法则（起点终点、分支数、实轴段、根之和、渐近线、分离点、虚轴交点、出射入射角）；
    \item 由根轨迹确定指定 $K^*$ 下的闭环极点、临界增益与稳定范围；
    \item 广义根轨迹：参数根轨迹与零度根轨迹；
    \item 用根轨迹定性分析稳定性、动态性能（主导极点）与稳态误差。
\end{enumerate}
\noindent\textbf{补充：}根轨迹法与时域法、频域法并列为经典控制理论的三大分析设计方法，1948 年由 W.\,R.\,伊文思（Evans）提出。它的全部结论都建立在\textbf{开环}零、极点已知这一前提上，这正是它能"不解高次方程而得到闭环极点"的原因。

\medskip

\section{根轨迹的基本概念}
\subsection{根轨迹的定义}
当系统某一参数（通常是开环增益 $K$ 或根轨迹增益 $K^*$）由 $0$ 连续增大到 $+\infty$ 时，闭环系统特征方程的根在 $s$ 平面上移动所描绘出来的轨迹，称为该系统的\textbf{根轨迹}。

两点说明：根轨迹是\textbf{闭环极点}的轨迹（有零、极点相消时被消去的极点不出现）；变化参数不限于 $K^*$，只要能把特征方程整理成 $A\,P(s)/Q(s)=-1$（$P,Q$ 为与 $A$ 无关的首一多项式），就可对 $A$ 画根轨迹，画多个参数得到根轨迹簇。

\paragraph{为什么要用根轨迹}
四阶以上特征方程无通用求根公式，参数每变一次都要重算，且只给出离散根值、看不出移动趋势。根轨迹把全部 $K^*$ 对应的闭环极点画在一张图上，三类性能都能直接读出；反过来，给定极点的容许区域（如 $\zeta\ge0.5$），可由图反求所需增益或校正装置参数。

\subsection{闭环零、极点与开环零、极点之间的关系}
设前向通路与反馈通路分别为
\[
G(s)=K_G^*\,\frac{\prod_{i=1}^{f}(s-z_i)}{\prod_{i=1}^{q}(s-p_i)},
\qquad
H(s)=K_H^*\,\frac{\prod_{j=1}^{l}(s-z_j)}{\prod_{j=1}^{h}(s-p_j)}
\]
则 $G(s)H(s)=K^*\prod_{i=1}^{m}(s-z_i)\big/\prod_{j=1}^{n}(s-p_j)$，其中 $K^*=K_G^*K_H^*$ 称为\textbf{开环根轨迹增益}，$m=f+l$，$n=q+h$。代入 $\Phi(s)=G/(1+GH)$ 得
\[
\Phi(s)=
\frac{K_G^*\displaystyle\prod_{i=1}^{f}(s-z_i)\prod_{j=1}^{h}(s-p_j)}
{\displaystyle\prod_{i=1}^{n}(s-p_i)+K^*\prod_{j=1}^{m}(s-z_j)}
\]
比较分子、分母，得三条结论：
\begin{enumerate}
    \item \textbf{闭环根轨迹增益}等于前向通路的根轨迹增益 $K_G^*$；单位反馈时即 $K^*$。
    \item \textbf{闭环零点}由前向通路零点与反馈通路极点组成；单位反馈时闭环零点就是开环零点。
    \item \textbf{闭环极点}与开环零点、开环极点以及 $K^*$ 都有关，不能由开环零、极点单独确定——这正是要用图解法找它的原因。
\end{enumerate}

\paragraph{易错点}非单位反馈时"闭环零点 $=$ 开环零点"是错的：反馈通路的极点会变成闭环零点。

\subsection{根轨迹与系统性能}
\begin{itemize}
    \item \textbf{稳定性}：全部闭环极点在左半平面则稳定。全部支路都在左半平面 $\Rightarrow$ 对一切 $K^*>0$ 稳定；与虚轴有交点 $\Rightarrow$ 交点处 $K^*$ 即\textbf{临界根轨迹增益}，稳定范围 $0<K^*<K^*_c$；有支路始终在右半平面 $\Rightarrow$ 任何 $K^*$ 都不稳定。
    \item \textbf{稳态性能}：原点处开环极点数 $v$ 决定型别，而 $v$ 就是根轨迹起点中落在原点的个数，图上可直接数出；图上 $K^*$ 换算成 $K$ 后即为该型别的静态误差系数。
    \item \textbf{动态性能}：实轴上的轨迹段对应实数极点，响应单调；进入复平面后对应共轭复极点，响应为衰减振荡。$K^*$ 增大使极点向虚轴移动，阻尼比减小、超调增大。
\end{itemize}

\subsection{根轨迹增益与开环增益的换算}
开环传递函数写成时间常数形式 $G(s)H(s)=K\prod(\tau_i s+1)\big/\left[s^{v}\prod(T_j s+1)\right]$，与根轨迹形式对照得
\[
K=K^*\,\frac{\displaystyle\prod_{i=1}^{m}|z_i|}{\displaystyle\prod_{j=1}^{n-v}|p_j|}
\]
$v$ 为原点处开环极点数；分母只对\textbf{非零}开环极点取积，$m=0$ 时分子取 1。即 $K$ 与 $K^*$ 只差一个比例常数。

\paragraph{易错点}图上标注的、用模值条件算出的都是 $K^*$，题目问的"开环增益"是 $K$；漏掉换算（或分母写成全部极点）会使答案整体差一个常数倍。

\medskip

\section{根轨迹方程与相角条件、模值条件}
\subsection{根轨迹方程的导出}
闭环特征方程 $1+G(s)H(s)=0$ 代入开环传递函数，得
\[
K^*\,\frac{\prod_{i=1}^{m}(s-z_i)}{\prod_{j=1}^{n}(s-p_j)}=-1
\]
称为\textbf{根轨迹方程}。它是向量方程，直接使用不便；注意到 $-1=1\cdot\mathrm{e}^{\mathrm{j}(2k+1)\pi}$，可拆成模与相角两个实数方程。

\subsection{相角条件与模值条件}
\textbf{相角条件}
\[
\sum_{i=1}^{m}\angle(s-z_i)-\sum_{j=1}^{n}\angle(s-p_j)=(2k+1)\pi,
\qquad k=0,\pm1,\pm2,\cdots
\]
用角度制即 $180^\circ+360^\circ k$，故常规根轨迹又称\textbf{$180^\circ$ 根轨迹}。

\textbf{模值条件}（也称幅值条件）
\[
K^*=\frac{\displaystyle\prod_{j=1}^{n}|s-p_j|}{\displaystyle\prod_{i=1}^{m}|s-z_i|}
\]
即 $K^*$ 等于各开环极点到 $s$ 的距离之积除以各开环零点到 $s$ 的距离之积。

\subsection{两个条件的地位}
\begin{itemize}
    \item \textbf{相角条件是充要条件}：满足它的点一定在根轨迹上，不满足的一定不在。画轨迹只用相角条件。
    \item \textbf{模值条件是必要条件}：对 $s$ 平面上\emph{任意}一点，总能找到一个 $K^*$ 使模值条件成立，但该点未必在轨迹上。模值条件只用来确定轨迹上某点对应的 $K^*$。
\end{itemize}
概括：\textbf{相角条件定形状，模值条件定参数。}

\subsection{用相角条件判定某点是否在根轨迹上}
\begin{enumerate}
    \item 标出全部开环零点 $z_i$ 与极点 $p_j$；
    \item 由每个零、极点向待判点 $s$ 作向量，模为 $|s-z_i|$ 或 $|s-p_j|$，相角为 $\angle(s-z_i)$ 或 $\angle(s-p_j)$（是\textbf{由零点/极点指向 $s$}）；
    \item 按相角条件求和，看是否为 $180^\circ$ 的奇数倍；满足者再用模值条件求该点的 $K^*$。
\end{enumerate}
共轭复极点（零点）对实轴上任意点的张角之和为 $360^\circ$ 的整数倍，对相角条件无贡献——这正是"实轴段法则"成立的原因。

\medskip

\section{绘制 $180^\circ$ 根轨迹的基本法则}
以下八条均以 $K^*$ 为变化参数、以负反馈系统为对象，画出的轨迹满足 $180^\circ+360^\circ k$ 的相角条件。

\subsection{法则 1：起点和终点}
\paragraph{结论}根轨迹起始于开环极点，终止于开环零点；若 $n>m$，有 $n-m$ 条分支终止于无穷远处。

\paragraph{条件与用途}对任何 $n\ge m$ 的系统成立，是画图的出发点，用于确定每条分支的起止。由 $K^*=\prod(s-p_j)/\prod(s-z_i)$ 可见：$K^*\to0$ 时右端为零只能靠分子为零，得 $s=p_j$；$K^*\to\infty$ 时右端发散只能靠分母为零或 $s\to\infty$。

\subsection{法则 2：分支数、对称性与连续性}
\paragraph{结论}分支数等于开环极点数 $n$（即系统阶数、闭环极点数）；轨迹连续，且关于实轴上下对称。

\paragraph{条件与用途}恒成立。用于清点分支、检查图是否画漏：$n$ 条分支从 $n$ 个极点出发，其中 $m$ 条终止于零点，其余 $n-m$ 条伸向无穷远。特征方程为实系数代数方程，复根必共轭成对，故对称于实轴；其根是系数的连续函数，故轨迹不断开。

\subsection{法则 3：实轴上的根轨迹}
\paragraph{结论}从实轴上\textbf{最右端}的开环零、极点起向左数，某段右侧（不含该点本身）的\textbf{实数}开环零、极点总数为\textbf{奇数}，则该段必是根轨迹；为偶数则不是。

\paragraph{条件与用途}只数实数零、极点，复零极点一律不计。用于快速确定实轴上的轨迹段。

\paragraph{易错点}把复极点也算进"个数"；或忘记从"最右端"起算，导致整条实轴段的判断整体反了。

\subsection{法则 4：根之和}
\paragraph{结论}当 $n-m\ge2$ 时，闭环极点之和为常数，等于开环极点之和：
\[
\sum_{i=1}^{n}\lambda_i=\sum_{j=1}^{n}p_j=-a_{n-1}=C
\]
即根之和与 $K^*$ 无关。

\paragraph{条件与用途}必须 $n-m\ge2$。用于：已知 $n-1$ 个闭环极点时直接求最后一个；判断"一部分根左移必有一部分右移、移动总量为零"；画完图后验算。理由：$D(s)=\prod(s-p_j)+K^*\prod(s-z_i)$ 中，$n-m\ge2$ 时 $K^*$ 项最高次不超过 $n-2$，不影响 $a_{n-1}$；$n-m=1$ 时 $K^*$ 项含 $s^{n-1}$，根之和随 $K^*$ 变化，本法则失效。

\subsection{法则 5：渐近线}
\paragraph{结论}当 $n>m$ 时，伸向无穷远的 $n-m$ 条分支沿 $n-m$ 条渐近线趋于无穷，其与实轴的交点及倾角为
\[
\sigma_a=\frac{\sum_{j=1}^{n}p_j-\sum_{i=1}^{m}z_i}{n-m},
\qquad
\varphi_a=\frac{(2k+1)\pi}{n-m},\quad k=0,1,\cdots,n-m-1
\]
角度取 $n-m$ 个不同值，相邻渐近线夹角 $2\pi/(n-m)$，整组关于实轴对称。

\paragraph{条件与用途}只在 $n>m$ 时存在；$n=m$ 时无渐近线。用于确定远端走向，是判断"是否进入右半平面"的关键。

\paragraph{易错点}把 $\sigma_a$ 分子写成"极点之和 $+$ 零点之和"（零点项是减）；或角度只取到 $\pi$ 而漏掉下半平面分支。

\subsection{法则 6：分离点与分离角}
\paragraph{结论}两条或两条以上分支在 $s$ 平面相遇又立即分开的点，称为\textbf{分离点}（由实轴进入复平面处）或\textbf{会合点}（由复平面回到实轴处），其坐标 $d$ 满足
\[
\sum_{j=1}^{n}\frac{1}{d-p_j}=\sum_{i=1}^{m}\frac{1}{d-z_i}
\]
等价于 $\left.\mathrm{d}D(s)/\mathrm{d}s\right|_{s=d}=0$（$d$ 为特征方程的重根）；开环无零点时右端取 $0$。若有 $l$ 条分支在此相遇，\textbf{分离角}为 $\theta_d=(2k+1)\pi/l$，$k=0,1,\cdots,l-1$；$l=2$ 时必为 $\pm90^\circ$。

\paragraph{条件与用途}用于确定分支离开/回到实轴的位置，以及"由单调变为振荡"的临界增益 $K^*_d$（在 $d$ 处由模值条件算得）。实轴上相邻两极点之间至少有一个分离点，相邻两零点之间至少有一个会合点，一极一零之间可能没有。

\paragraph{易错点}分离点方程的解只是\textbf{必要条件}，须回到相角条件检验是否真在轨迹上，不在的舍去；$d$ 处由模值条件算出的 $K^*$ 必须为正。

\subsection{法则 7：与虚轴的交点}
\paragraph{结论}交点处的 $K^*$ 即\textbf{临界根轨迹增益}，交点纵坐标即临界振荡频率 $\omega_c$。两种求法：

\emph{方法一（代入 $s=\mathrm{j}\omega$）}：令 $D(\mathrm{j}\omega)=0$，分离实部与虚部
\[
\mathrm{Re}\left[D(\mathrm{j}\omega)\right]=0,
\qquad
\mathrm{Im}\left[D(\mathrm{j}\omega)\right]=0
\]
联立解出 $\omega$ 与 $K^*$（$\omega=0$ 的解对应轨迹穿过原点，不要漏）。

\emph{方法二（劳斯判据）}：令劳斯表出现\textbf{全零行}（临界稳定），由全零行上一行的辅助方程解出 $\omega_c$，由全零行条件解出 $K^*_c$。

\paragraph{条件与用途}用于确定稳定范围 $0<K^*<K^*_c$、临界振荡频率，判断轨迹是否进入右半平面。

\paragraph{易错点}漏掉 $\omega=0$ 的解；或把 $K^*_c$ 直接当开环增益 $K$ 报出。

\subsection{法则 8：出射角与入射角}
\paragraph{结论}轨迹离开复数开环极点 $p_k$ 的方向称\textbf{出射角}（起始角），进入复数开环零点 $z_k$ 的方向称\textbf{入射角}（终止角）。二者都由相角条件解出未知角得到：把 $s$ 取在待求的复数极点（零点）上，除所求角外其余各项已知，解方程即得。写成显式：
\[
\theta_{p_k}=(2k+1)\pi+\sum_{i=1}^{m}\angle(p_k-z_i)-\sum_{\substack{j=1\\ j\neq k}}^{n}\angle(p_k-p_j)
\]
\[
\varphi_{z_k}=(2k+1)\pi-\sum_{\substack{i=1\\ i\neq k}}^{m}\angle(z_k-z_i)+\sum_{j=1}^{n}\angle(z_k-p_j)
\]

\paragraph{条件与用途}只对\textbf{复数}开环零、极点才需要计算；实轴上的极点出射角只能是 $0^\circ$ 或 $180^\circ$。用于确定分支离开复极点、进入复零点时的切线方向，否则复平面段画不出来。

\paragraph{易错点}符号方向弄反：出射角中零点项取\textbf{加}、其余极点项取\textbf{减}；入射角中其余零点项取\textbf{减}、极点项取\textbf{加}。

\subsection{绘制概略根轨迹的步骤}
\begin{enumerate}
    \item \textbf{化标准形}：写成 $G(s)H(s)=K^*\prod(s-z_i)/\prod(s-p_j)$，标出 $n$、$m$、全部零极点，算出 $K$ 与 $K^*$ 的换算系数。
    \item \textbf{定端点、数分支}（法则 1、2）：标出 $n$ 个起点、$m$ 个终点。
    \item \textbf{定实轴段}（法则 3）：从最右端零极点起向左逐段判定。
    \item \textbf{定渐近线}（法则 5，仅 $n>m$）：算 $\sigma_a$ 与 $\varphi_a$，画出 $n-m$ 条射线。
    \item \textbf{定分离点}（法则 6）：解分离点方程，舍去不在轨迹上的根，算出 $d$ 处的 $K^*_d$ 与分离角。
    \item \textbf{定虚轴交点}（法则 7）：求 $\omega_c$ 与 $K^*_c$。
    \item \textbf{定出射角、入射角}（法则 8）。
    \item \textbf{校验并连线}：用根之和（法则 4）校验极点和；必要时补算几个点，光滑连线并沿 $K^*$ 增大方向画箭头。
\end{enumerate}

\paragraph{指定 $K^*$ 下闭环极点的确定}先用模值条件在实轴段上试探求出实数闭环极点，再用综合除法降阶解出其余极点。若给的是阻尼比 $\zeta$（等阻尼线，过原点的射线）或自然频率 $\omega_n$（以原点为心的圆），则先由根轨迹与该线/圆的交点定出极点位置，再用模值条件反求 $K^*$。

\subsection{八条法则一览}
{\renewcommand{\arraystretch}{1.3}\setlength{\tabcolsep}{6pt}
\begin{table}[htbp]
\centering
\caption{$180^\circ$ 常规根轨迹绘制法则一览}
\begin{tabular}{@{}p{1.0cm}p{5.7cm}p{3.2cm}p{4.0cm}@{}}
\toprule
法则 & 结论 & 条件 & 用途 \\
\midrule
1 & 起于开环极点，终于开环零点；$n-m$ 条终于无穷远 & $n\ge m$ & 定端点与分支去向 \\
2 & 分支数 $=n$；连续；对称于实轴 & 恒成立 & 清点分支、查漏 \\
3 & 右侧实数零极点个数为奇数的实轴段是轨迹 & 只数实数零极点 & 定实轴段 \\
4 & $\sum\lambda_i=\sum p_j=C$ & $n-m\ge2$ & 求末极点、验算 \\
5 & $\sigma_a=\frac{\sum p-\sum z}{n-m}$，\newline $\varphi_a=\frac{(2k+1)\pi}{n-m}$ & $n>m$ & 定远端走向 \\
6 & $\sum\frac{1}{d-p_j}=\sum\frac{1}{d-z_i}$，\newline $\theta_d=\frac{(2k+1)\pi}{l}$ & 有分支相遇 & 定分离/会合点 \\
7 & $\mathrm{Re}D(\mathrm{j}\omega)=\mathrm{Im}D(\mathrm{j}\omega)=0$\newline 或劳斯表全零行 & 轨迹穿越虚轴 & 定临界增益与稳定范围 \\
8 & 由相角条件解出未知角 & 复数零极点 & 定复平面段切向 \\
\bottomrule
\end{tabular}
\end{table}}

\medskip

\section{复平面根轨迹的两个几何定理}
\paragraph{定理 1（圆弧定理）}若系统有 2 个开环极点、1 个开环零点，且在复平面内存在根轨迹，则该复平面根轨迹一定是以该零点为圆心的圆弧——此时相角条件可化为 $|s-z_1|=$ 常数，半径由分离点到 $z_1$ 的距离决定。

\paragraph{定理 2（对称性定理）}若开环零、极点均为偶数个，且关于一条平行于虚轴的直线左右对称分布，则根轨迹一定关于该直线左右对称。两定理都用于判断复平面段的形状，减少补点计算。

\medskip

\section{广义根轨迹}
\subsection{含义}
除根轨迹增益 $K^*$ 变化以外的其他情形下的根轨迹，统称\textbf{广义根轨迹}，包括参数根轨迹、零度根轨迹、以及 $m>n$ 时的根轨迹等。相对地，负反馈系统中 $K^*$ 由 $0\to\infty$ 变化时的根轨迹称为\textbf{常规根轨迹}（$180^\circ$ 根轨迹）。

\subsection{参数根轨迹}
以非开环增益的参数 $A$（如时间常数 $T$、测速反馈系数、附加零点位置）为可变参数绘制的根轨迹，称为\textbf{参数根轨迹}。

\paragraph{构造等效开环传递函数}把闭环特征方程整理、分离出含 $A$ 与不含 $A$ 的项，写成
\[
A\,\frac{P(s)}{Q(s)}=-1
\]
$P(s)$、$Q(s)$ 为与 $A$ 无关的首一多项式。于是可把系统看成"开环传递函数为 $G_1(s)H_1(s)=A\,P(s)/Q(s)$ 的单位反馈系统"，对 $A$ 按八条法则作图：$A$ 相当于 $K^*$，$P(s)$ 的根是等效开环零点，$Q(s)$ 的根是等效开环极点。

\paragraph{关键提醒（易错）}"等效"只在\textbf{闭环极点相同}这一点上成立，闭环零点一般不同。故分析性能时：\textbf{极点}可取自参数根轨迹，\textbf{零点}必须回到原闭环系统去取；算稳态误差也必须用原系统的开环增益，不能把 $A$ 当 $K^*$ 代入。

\subsection{零度根轨迹}
\paragraph{来源}零度根轨迹的相角条件为 $0^\circ+2k\pi$，出现在两种场合：非最小相位系统（开环传递函数含 $s$ 最高次幂系数为负的因子）；控制系统中含有\textbf{正反馈}内回路。

\paragraph{根轨迹方程}正反馈内回路闭环传递函数为 $G_2/(1-G_2H_2)$，特征方程 $1-G_2H_2=0$，即 $G_2(s)H_2(s)=+1$，拆成
\[
\sum_{i=1}^{m}\angle(s-z_i)-\sum_{j=1}^{n}\angle(s-p_j)=2k\pi
\quad\text{（相角条件）}
\]
\[
K^*=\frac{\displaystyle\prod_{j=1}^{n}|s-p_j|}{\displaystyle\prod_{i=1}^{m}|s-z_i|}
\quad\text{（模值条件，与 $180^\circ$ 完全相同）}
\]

\paragraph{需要修改的三条法则}模值条件没变，故与模值无关的法则（1、2、4、6、7）照用；与相角有关的要改：
\begin{enumerate}
    \item \textbf{实轴段}（对应法则 3）：右侧实数开环零、极点个数之和为\textbf{偶数}的区段才是根轨迹（原为奇数）。
    \item \textbf{渐近线倾角}（对应法则 5）：$\varphi_a=2k\pi/(n-m)$，$k=0,1,\cdots,n-m-1$；交点 $\sigma_a$ 公式不变。
    \item \textbf{出射角、入射角}（对应法则 8）：把公式中的 $(2k+1)\pi$ 换成 $2k\pi$。
\end{enumerate}

\paragraph{易错点}只改相角条件却忘了同步修改实轴段、渐近线倾角、出射角三处，图形会整体错位。口诀："\textbf{模值不变，相角全换}"。

\medskip

\section{用根轨迹分析系统性能}
\subsection{分析步骤}
\begin{enumerate}
    \item 按八条法则画出根轨迹（参数根轨迹先做等效变换，零度根轨迹先改相角相关法则）；
    \item \textbf{稳定性}：看是否有分支进入右半平面，由与虚轴交点定临界增益与稳定范围；
    \item \textbf{动态性能}：由给定 $K^*$ 定出闭环极点，判断运动形式，找出主导极点，用二阶指标公式估算 $\sigma\%$ 与 $t_s$；
    \item \textbf{稳态性能}：由原点处开环极点数定型别，由 $K^*$ 换算 $K$，算稳态误差；
    \item 指标不满足时先调 $K^*$，仍不满足则改变结构或加校正装置（第 6 章）。
\end{enumerate}

\subsection{稳定性分析}
\begin{itemize}
    \item 闭环极点全部在左半平面 $\Longleftrightarrow$ 系统稳定。\textbf{稳定性只取决于闭环极点的位置，与闭环零点无关}。
    \item 与虚轴的交点给出临界根轨迹增益 $K^*_c$ 与临界振荡频率 $\omega_c$；若 $K^*=0$ 时轨迹就在左半平面，则稳定范围为 $0<K^*<K^*_c$。
    \item 某条分支随 $K^*$ 增大穿过虚轴进入右半平面，则 $K^*>K^*_c$ 时不稳定，通常表现为振幅不断增大的振荡。
    \item $K^*_c$ 须换算成开环增益 $K_c$；题目问"$K$ 的范围"时不能直接写 $K^*$ 的范围。
\end{itemize}

\subsection{动态性能与主导极点}
\paragraph{由分离点划分运动形式}设分离点处增益为 $K^*_d$：$0<K^*<K^*_d$ 时全部闭环极点为实数，单位阶跃响应单调上升（过阻尼或临界阻尼），无超调；$K^*>K^*_d$ 时出现一对共轭复极点，响应为衰减振荡（欠阻尼），超调随 $K^*$ 增大而增大。

\paragraph{极点位置与动态指标的对应}设主导极点为 $\lambda_{1,2}=-\zeta\omega_n\pm\mathrm{j}\omega_n\sqrt{1-\zeta^2}$，则
\[
\zeta=\cos\theta\ \ (\theta\ \text{为极点向量与负实轴的夹角}),\qquad
\omega_n=|\lambda|,\qquad
t_s\approx\frac{3\sim4}{\zeta\omega_n}
\]
于是 $s$ 平面上有三族等值线：过原点的\textbf{等阻尼线}（射线，$\zeta$ 相同）、以原点为心的\textbf{等 $\omega_n$ 圆}、平行于虚轴的\textbf{等 $t_s$ 线}（$|\mathrm{Re}\,\lambda|$ 相同）。根轨迹与这些线的交点即满足相应指标的极点位置，再由模值条件反求 $K^*$。

\paragraph{主导极点的定义}在全部闭环极点中，\textbf{最靠近虚轴、且其附近没有闭环零点}、对整个时间响应起主要作用的那些闭环极点，称为\textbf{主导极点}；可以是实数极点，也可以是一对共轭复极点。"附近无闭环零点"不可省：响应分量的消逝速度还取决于该极点处的留数，若某极点旁紧挨一个闭环零点，该分量被大幅削弱，即使靠近虚轴也未必是主导极点。

\paragraph{偶极子与主导极点法}若一对闭环零、极点相距很近（经验判据：二者间距比它们本身的模值小一个数量级），则称它们构成\textbf{偶极子}；复数偶极子必共轭出现，远离原点的可忽略，接近原点的必须考虑。\textbf{主导极点法}即选留最靠近虚轴而又不十分靠近闭环零点的一两个（或两三个）极点，略去不十分接近原点的偶极子以及比主导极点距虚轴远 3～6 倍以上的其他闭环零、极点，把高阶系统简化为低阶系统。注意：略去零极点后闭环增益常会改变，必须重新核算 $\Phi(0)$，否则会得出"稳态误差不为零"的错误结论；输入信号的极点不在选择范围之内。

\paragraph{闭环零、极点对动态性能的影响（定性规律）}
\begin{enumerate}
    \item \textbf{稳定性与运动形式}：稳定性只与闭环极点位置有关，与零点无关；无零点且极点全为实数时响应单调，极点全为复数时响应一般振荡。
    \item \textbf{超调量}：主要取决于闭环复数主导极点的衰减率 $\zeta$，即 $\sigma\%=\mathrm{e}^{-\pi\zeta/\sqrt{1-\zeta^2}}$，并与其他零、极点接近原点的程度有关。
    \item \textbf{调节时间}：主要取决于最靠近虚轴的闭环复数极点的实部绝对值 $\zeta\omega_n$；若某实数极点距虚轴最近且附近无实数零点，则 $t_s$ 主要由该实数极点的模值决定。
    \item \textbf{实数零点}：减小系统阻尼，使峰值时间提前、超调增大；\textbf{实数极点}：增大系统阻尼，使峰值时间滞后、超调减小。二者作用都随其接近原点的程度而加强。
\end{enumerate}

\paragraph{易错点}把"零点使超调增大"记成"减小"、把"极点使峰值时间滞后"记成"提前"——两者方向恰好相反。记忆方法：零点相当于微分（加快、加剧），极点相当于惯性滤波（滞后、平滑）。

\subsection{稳态误差分析}
\begin{itemize}
    \item \textbf{型别}由原点处开环极点数 $v$ 决定，而 $v$ 就是根轨迹起点中落在原点的个数，图上可直接数出。
    \item 0 型：阶跃输入下 $e_{ss}=1/(1+K)$；I 型：阶跃无差，斜坡输入下 $e_{ss}=1/K$（$K=K_v$）；II 型：斜坡无差，等加速输入下 $e_{ss}=1/K$（$K=K_a$）。
    \item 流程：由该点位置用模值条件求 $K^*$ $\to$ 按换算公式求 $K$ $\to$ 由型别选公式算 $e_{ss}$。
    \item 若用参数根轨迹，图上参数是 $A$ 而不是 $K^*$，必须先由闭环极点反算出原系统的 $K$，再算稳态误差。
\end{itemize}

\medskip

\section{本章易错点汇总}
{\renewcommand{\arraystretch}{1.35}\setlength{\tabcolsep}{6pt}
\begin{table}[htbp]
\centering
\caption{根轨迹法常见错误与正确做法}
\begin{tabular}{@{}p{4.3cm}p{7.0cm}p{2.5cm}@{}}
\toprule
易错点 & 正确做法 & 涉及 \\
\midrule
用模值条件判断某点是否在轨迹上 & 模值条件对任意点都能凑出 $K^*$；判轨迹只能用相角条件 & 条件地位 \\
把开环零点当闭环零点 & 非单位反馈时闭环零点 $=$ 前向零点 $+$ 反馈极点 & 零极点关系 \\
$K$ 与 $K^*$ 混用 & $K=K^*\prod|z_i|/\prod|p_j|$，分母只取非零极点 & 换算 \\
实轴段把复极点也数进去 & 只数实轴上的实数零、极点，且从最右端算起 & 3 \\
根之和法则无条件使用 & 仅 $n-m\ge2$ 成立；$n-m=1$ 时根之和随 $K^*$ 变 & 4 \\
$n=m$ 时也画渐近线 & $n>m$ 才有渐近线，条数 $n-m$ & 5 \\
分离点方程的解全部保留 & 解出的 $d$ 须验证在轨迹上，且 $K^*_d>0$ & 6 \\
求虚轴交点漏掉 $\omega=0$ & 实部、虚部方程联立，$\omega=0$ 的解也要检验 & 7 \\
出射角符号弄反 & 极点：零点项取加、其余极点项取减 & 8 \\
零度根轨迹只改相角条件 & 实轴段（偶）、渐近线倾角、出射角三处都要改 & 广义 \\
用等效开环分析性能 & 极点可用，零点与稳态误差必须回原系统 & 参数根轨迹 \\
略去零极点后不校验终值 & 重新核算 $\Phi(0)$，防止稳态误差算错 & 性能分析 \\
\bottomrule
\end{tabular}
\end{table}}

% ===== 恢复全局默认间距（避免影响后续章节）=====
\setlength{\parskip}{0pt plus 1pt}
\setlist{leftmargin=2em, itemsep=0.2em, parsep=0pt}
\renewcommand{\arraystretch}{1}
\setlength{\tabcolsep}{6pt}
% ============================================
